\documentclass[11pt,a4paper]{article}
\usepackage[utf8]{utf8}
\usepackage[margin=1in]{geometry}
\usepackage{amsmath,amssymb}
\usepackage{enumitem}


\begin{document}


\begin{center}
    {\Large   \textbf{B.Sc. (Hons.) Semester III (NEP) Mid-Semester Examination 2025--26}}\\[4pt]
    {\large   \textbf{DST-Centre for Interdisciplinary Mathematical Sciences (CIMS)}}\\[4pt]
   \textbf{Paper Code: CIMSE31 : Python Programming for Data Analytics}\\[6pt]
   \textbf{Time: 1 Hour} \hfill   \textbf{Max. Marks: 20}
\end{center}
\hr




\noindent   \textbf{Note:} Attempt any FOUR questions. Each question carries equal marks (5 Marks each).


\begin{enumerate}[label=   \textbf{\arabic*.}]
    \item What do you understand by Object-Oriented Programming (OOP)? Explain the core features of Python that make it unique and beneficial compared to other programming languages. \hfill [5]
    \item Explain the key differences among Tuple, List, and Dictionary in Python with suitable code examples. Highlight their mutability and primary use cases. \hfill [5]
    \item Write a Python script to perform basic data manipulation using Pandas. Discuss the creation of Pandas Series and DataFrame objects. \hfill [5]
    \item Explain conditional statements and loop control mechanisms (   \texttt{for},   \texttt{while},   \texttt{break},   \texttt{continue}) in Python with appropriate examples. \hfill [5]
    \item Discuss the role of Python libraries such as NumPy and Matplotlib in Data Analytics with relevant snippets. \hfill [5]
\end{enumerate}

\end{document}
